\lecture{13}{Wed 28 Sep 2022 11:31}{}

\begin{lemma}
	$x_n \to a$ in $\mathbb{R}^p$ if and only of $\|x_n - a\| \to 0$ in $\mathbb{R}^p$.
\end{lemma}
\begin{lemma}
	Let $\alpha_n>0, \alpha_n \to 0$. If $\|x_n - a\|\le C \alpha_n$ for all $n \ge K_0$, then $x_n \to  a$.
\end{lemma}
\begin{proof}
	If $\alpha_n > \varepsilon$ for $n \ge K_\varepsilon$, then $\|x_n - a\|< C\varepsilon$ for $n \ge K_\varepsilon$. We let $\varepsilon' =  C \varepsilon $.
\end{proof}

\begin{example}
	\begin{enumerate}
		\item $\frac{1}{n} \to  0$ as $n \to \infty $
		\item Let $0<|b|<1$. Then $b^n \to 0$ as $n\to 0$.
			\begin{proof}
				Represent $|b| = \frac{1}{1+\alpha}, \alpha>0$. By Bernoulli, $(1+\alpha)^n \ge 1+n \alpha$. Thus $|b|^n \le \frac{1}{1+n \alpha}\le \frac{1}{n \alpha }$. 

				But then $|b|^n \le \frac{1}{n \alpha}< \varepsilon$ if $n>\frac{1}{\alpha \varepsilon}$. Since $|b^n - 0| = |b|^n < \varepsilon$ if $n > \frac{1}{\alpha\varepsilon}$, so $b^n \to 0$.
			\end{proof}
		\item Let $a>0$. Then $a^{\frac{1}{n}}\to 1$ as $n\to \infty $.
			\begin{proof}
				Assume  $a>1$. Let $d_n = a^{\frac{1}{n}}-1> 0$ since $a>1$. 
				Now $a = (1+d_n)^n\ge 1+ n d_n$.
				Hence $d_n \le \frac{a-1 }{n}<\varepsilon$ if $n\ge \frac{a-1}{\varepsilon}$.
				Hence $|d_n|<\varepsilon$ if $n \ge K_\varepsilon$, so $d_n \to  0$.

			Now assume $a \le 1$. We prove later that $x_n \to b$ implies $\frac{1}{x_n}\to \frac{1}{b}$. Then we use the previous argument to see $\frac{1}{a^{\frac{1}{n}}}\to 1$, thus $a^{\frac{1}{n}}\to 1$.
			\end{proof}
		\item $n^{\frac{1}{n}}\to 1$. 
			\begin{proof}
				Let $c_n = n^{\frac{1}{n}}-1\ge 0$.
				Now $(1+c_n)^n = n$. We use 2 terms in the Binomial Expansion to bound $n$ from below:
				\[
					n \ge 1+n c_n + \frac{n(n-1)}{2}c_n^2
				.\] 
				We have used the 
				\begin{theorem}[Newton's Binomial Theorem]
					\[
						(1+\alpha)^n = 1 + n \alpha + \ldots + (n,k) \alpha^k + \ldots + \alpha^n
					.\] 
					where \[
						(n,k) = \frac{n!}{k! (n-k)!}
					.\] 
				\end{theorem}
				Now $c_n^2 \le \frac{2}{n-1} \to 0$. Hence $c_n \to 0$.
			\end{proof}
			
			
			
	\end{enumerate}
\end{example}

\subsection{Subsequences}

\begin{definition}[Subsequence]
	Let $X$ be a sequence and $r_1 < r_2 <\ldots<r_n$ be a sequence of natural numbers.
	Then $X' = (x_{r_1}, x_{r_2}, \ldots,x_{r_n},\ldots)$ is called a \textit{subsequence} of $X$.
\end{definition}

\begin{example}
	\begin{enumerate}
		\item $(x_{2n})$ is a subsequence of $(x_n)$.
		\item Tails: $(x_n)_{n\ge k}$ is called a \textit{tail} of $X$. It is a subsequence of $(x_n)$.
	\end{enumerate}
\end{example}

\begin{theorem}
	If $X$ is a sequence in $\mathbb{R}^p$, and $\lim X = a$, then any subsequence $X'$ of $X$ also converges to $a$.
\end{theorem}
\begin{proof}
	For any $\varepsilon>0$, there exists $K_\varepsilon$ such that $\|x_n - a\|<\varepsilon$ if $n > K_\varepsilon$. 
	Let $X' = (x_{r_n})$. Note $r_n \ge n$. Hence \[
		\|x_{r_n} - a\|<\varepsilon \quad \text{for } n\ge K_\varepsilon
	.\] 
\end{proof}

\begin{example}
	$(-1)^n$ is divergent.
	\begin{proof}
		$x_{2n} = (-1)^{2n} = 1 \to  1$. But $x_{2n+1} = (-1)^{2n+1}= -1 \to -1$. If $X \to a$ then $x_{2n} \to a$ and $x_{2n+1}\to a$, so $-1 = 1$, a contradiction.
	\end{proof}
\end{example}

\begin{proposition}
	If $X \not \to a$, then there exists a subsequence of $X$ entirely out of some neighborhod of $a$.
\end{proposition}
